\documentclass[a4paper,12pt]{article}
\usepackage[T2A]{fontenc}
\usepackage[utf8]{inputenc}
\usepackage[russian]{babel}
\usepackage{amsmath,amssymb,mathtools}
\usepackage{geometry}
\geometry{left=24mm,right=24mm,top=23mm,bottom=24mm}
\usepackage{microtype}
\usepackage{icomma}
\usepackage{parskip}
\usepackage{setspace}
\usepackage{enumitem}
\usepackage{booktabs,tabularx}
\usepackage{xcolor}
\usepackage{tikz}
\usetikzlibrary{arrows.meta,positioning,shapes.geometric,calc}
\usepackage{fancyhdr}
\usepackage[hidelinks]{hyperref}
\definecolor{accent}{HTML}{008080}
\definecolor{darkaccent}{HTML}{005457}
\definecolor{textgray}{HTML}{393939}
\definecolor{muted}{HTML}{6B7778}
\color{textgray}
\onehalfspacing
\setlength{\parindent}{0pt}
\setlength{\parskip}{6pt}
\setlist[itemize]{leftmargin=*,itemsep=2pt,topsep=4pt}
\setlist[enumerate]{leftmargin=*,itemsep=6pt,topsep=4pt}
\pagestyle{fancy}
\fancyhf{}
\fancyhead[L]{\small Деление дробей, уравнения и текстовые задачи}
\fancyhead[R]{\small Надежда Клименко}
\fancyfoot[C]{\small\thepage}
\renewcommand{\headrulewidth}{0.4pt}
\renewcommand{\footrulewidth}{0pt}
\setlength{\headheight}{15pt}
\makeatletter
\def\headrule{{\color{accent}\hrule\@height\headrulewidth\@width\headwidth \vskip-\headrulewidth}}
\makeatother
\newcommand{\sect}[1]{\section*{#1}\vspace{-3pt}\noindent\textcolor{accent}{\rule{\linewidth}{0.5pt}}\vspace{3pt}}
\newcommand{\key}[1]{\textcolor{darkaccent}{\textbf{#1}}}
\newcommand{\ex}[1]{\medskip\noindent\textcolor{darkaccent}{\textbf{#1}}\par\vspace{-2mm}}
\newcommand{\mfrac}[3]{#1\,\frac{#2}{#3}}
\begin{document}
\thispagestyle{empty}
\begin{center}
\vspace*{28mm}
{\Large\textcolor{darkaccent}{МАТЕМАТИКА}}\par
\vspace{21mm}
{\LARGE Чек-лист}\par
\vspace{15mm}
{\LARGE\textcolor{accent}{Деление дробей, уравнения}}\par
\vspace{2mm}
{\LARGE\textcolor{accent}{и сложные текстовые задачи}}\par
\vspace{12mm}
{\large Повторение и углубление школьной математики}\par
\vfill
{\large Надежда Клименко}\par
\vspace{12mm}
{\large 27 октября 2026 г.}\par
\vspace{18mm}
{\normalsize Лёвин Артём Александрович}\par
{\normalsize эксклюзивно для Надежды Клименко}\par
\vspace{15mm}
\begin{tikzpicture}[scale=1]
\draw[accent,line width=0.85pt] (0,0) .. controls (1.2,0.65) and (2.0,-0.52) .. (3.1,0)
 .. controls (4.2,0.48) and (5.0,-0.48) .. (6.1,0)
 .. controls (7.2,0.55) and (7.9,-0.38) .. (9.2,0);
\fill[darkaccent] (4.6,0) circle (1.4pt);
\end{tikzpicture}
\end{center}
\clearpage

\sect{1. Как выполнять действия с дробями}
\key{Порядок действий:} скобки и степени $\rightarrow$ умножение и деление $\rightarrow$ сложение и вычитание. Действия одной ступени выполняйте слева направо.

\key{Смешанное число} сначала удобно перевести в неправильную дробь:
\[
1\frac15=\frac65,\qquad
\left(1\frac15\right)^2=\frac65\cdot\frac65=\frac{36}{25}=1{,}44.
\]

\key{Деление дробей:} чтобы разделить на дробь, умножьте на обратную ей дробь:
\[
\frac ab:\frac cd=\frac ab\cdot\frac dc,
\quad b\ne0,\ c\ne0,\ d\ne0.
\]
Если обе дроби десятичные, можно умножить делимое и делитель на одну и ту же степень десяти. Частное от этого не изменится:
\[
0{,}36:0{,}03=\frac{0{,}36}{0{,}03}
=\frac{36}{3}=12.
\]

\ex{Пример 1. Несколько действий в одном выражении}
\[
\left(\left(1\frac15\right)^2-1{,}08\right):0{,}03.
\]
\[
\begin{aligned}
1)&\quad \left(1\frac15\right)^2=1{,}44;\\
2)&\quad 1{,}44-1{,}08=0{,}36;\\
3)&\quad 0{,}36:0{,}03=12.
\end{aligned}
\]
\key{Проверка:} $12\cdot0{,}03=0{,}36$. Особое внимание второму действию: разность $1{,}44-1{,}08$ равна именно $0{,}36$.

\ex{Пример 2. Можно сочетать разные формы записи}
\[
9:0{,}18-37\frac12\cdot0{,}64
=50-37{,}5\cdot0{,}64
=50-24=26.
\]
\key{Выбор записи:} десятичные дроби удобны при делении на $0{,}18$ и умножении на $0{,}64$; обыкновенные дроби часто упрощают сокращение.

\sect{2. Обратные действия в уравнениях}
\begin{center}
\begin{tabular}{@{}p{0.28\linewidth}p{0.33\linewidth}p{0.33\linewidth}@{}}
\toprule
\textbf{Вид} & \textbf{Что ищем} & \textbf{Правило} \\
\midrule
$a\cdot x=c$ & неизвестный множитель & $x=c:a$, $a\ne0$ \\
$a:x=c$ & неизвестный делитель & $x=a:c$, $a,c\ne0$ \\
$a+x=c$ & неизвестное слагаемое & $x=c-a$ \\
$a-x=c$ & неизвестное вычитаемое & $x=a-c$ \\
\bottomrule
\end{tabular}
\end{center}
При записи $a:x=c$ подразумевается $x\ne0$. После решения полезно подставить найденное число в исходное уравнение.

\ex{Пример 3. Неизвестный множитель}
\[
\frac27x=1\frac17
\quad\Longrightarrow\quad
x=\frac87:\frac27=\frac87\cdot\frac72=4.
\]

\ex{Пример 4. Прежде чем делить, упростите правую часть}
\[
\frac35n=2\frac7{10}-\frac35
=\frac{27}{10}-\frac6{10}=\frac{21}{10}.
\]
\[
n=\frac{21}{10}:\frac35=\frac{21}{10}\cdot\frac53=\frac72=3\frac12.
\]

\ex{Пример 5. Сначала перенесите известное слагаемое}
\[
\frac49b+\frac37=1
\quad\Longrightarrow\quad
\frac49b=1-\frac37=\frac47
\quad\Longrightarrow\quad
b=\frac47:\frac49=\frac97.
\]
\clearpage

\sect{3. Подобные слагаемые}
\key{Подобные слагаемые} имеют одинаковую буквенную часть. Их можно складывать, складывая только числовые коэффициенты:
\[
ax+bx=(a+b)x,\qquad ax-bx=(a-b)x.
\]
Если перед буквой коэффициент не написан, он равен единице: $x=1\cdot x$.

\ex{Пример 6. Два подобных слагаемых}
\[
\frac19x+\frac49x=3\frac1{18}
\quad\Longrightarrow\quad
\frac59x=\frac{55}{18}.
\]
\[
x=\frac{55}{18}:\frac59
=\frac{55}{18}\cdot\frac95
=\frac{11}{2}=5\frac12.
\]

\ex{Пример 7. Три подобных слагаемых}
\[
\frac58x+x-\frac34x=1\frac34.
\]
\[
\left(\frac58+\frac88-\frac68\right)x=\frac74
\quad\Longrightarrow\quad
\frac78x=\frac74
\quad\Longrightarrow\quad
x=\frac74\cdot\frac87=2.
\]
\key{Частая ошибка:} после сведения подобных слагаемых остановиться на уравнении вида $\frac78x=\frac74$. Неизвестное ещё нужно найти делением.

\sect{4. Уравнения со скобками: «раскрываем» по действиям}
При работе с произведением вида $a\cdot(\ldots)=c$ удобно считать скобку одним неизвестным множителем. Сначала найдите её значение, затем последовательно находите число внутри скобки.

\ex{Пример 8. Составное уравнение}
\[
1\frac23\left(\frac13n+\frac37\right)=2\frac14.
\]
\textbf{Шаг 1.} Скобка --- неизвестный множитель:
\[
\frac13n+\frac37=\frac94:\frac53
=\frac94\cdot\frac35=\frac{27}{20}.
\]
\textbf{Шаг 2.} Найдите неизвестное слагаемое:
\[
\frac13n=\frac{27}{20}-\frac37
=\frac{189-60}{140}=\frac{129}{140}.
\]
\textbf{Шаг 3.} Найдите множитель $n$:
\[
n=\frac{129}{140}:\frac13
=\frac{387}{140}=2\frac{107}{140}.
\]
\key{Проверка:} внутри скобок $\frac{129}{140}+\frac37=\frac{27}{20}$, а $\frac53\cdot\frac{27}{20}=\frac94$.

\textbf{Важный знак:} из $a-b=c$ следует $a=c+b$. При переносе вычитаемого в другую часть уравнения оно становится слагаемым.

\ex{Пример 9. Неизвестный делитель}
\[
1\frac57:x=\frac67:2.
\]
\[
\frac67:2=\frac37,
\qquad \frac{12}{7}:x=\frac37,
\qquad x=\frac{12}{7}:\frac37=4.
\]
\clearpage

\thispagestyle{fancy}
\vspace*{8mm}
\begin{center}
{\Large\textcolor{darkaccent}{Алгоритм: как решать составное уравнение}}\par
\vspace{15mm}
\end{center}
\begin{center}
\tikzset{
 start/.style={ellipse,draw=accent,line width=.8pt,align=center,minimum width=43mm,minimum height=12mm},
 act/.style={rectangle,rounded corners=2mm,draw=accent,line width=.7pt,align=center,text width=92mm,minimum height=16mm,inner sep=5pt},
 decision/.style={diamond,aspect=1.8,draw=darkaccent,line width=.7pt,align=center,text width=27mm,inner sep=2pt},
 link/.style={-{Stealth[length=2.8mm]},line width=.7pt,draw=darkaccent}
}
\begin{tikzpicture}[node distance=12mm]
\node[start] (a) {Начало};
\node[act,below=of a] (b) {Выполните числовые действия,\\ приведите подобные слагаемые};
\node[decision,below=17mm of b] (c) {Получили $x=c$?};
\node[act,below=20mm of c] (d) {Найдите последнее действие\\ с неизвестным или скобкой};
\node[act,below=of d] (e) {Примените обратное действие\\ к обеим частям уравнения};
\node[start,right=7mm of c,minimum width=37mm] (f) {Проверка\ подстановкой};
\draw[link] (a)--(b);
\draw[link] (b)--(c);
\draw[link] (c)--node[right,font=\small]{нет}(d);
\draw[link] (d)--(e);
\draw[link] (e.west)--++(-13mm,0)|-([xshift=-13mm]b.west)--(b.west);
\draw[link] (c.east)--node[above,font=\small]{да}(f.west);
\end{tikzpicture}
\end{center}
\vfill
\begin{center}
\small \textcolor{muted}{Повторяйте шаги до тех пор, пока неизвестная не останется одна.}
\end{center}
\clearpage

\sect{5. Как переводить текст задачи в уравнение}
\begin{center}
\begin{tabular}{@{}p{0.35\linewidth}p{0.55\linewidth}@{}}
\toprule
\textbf{Формулировка} & \textbf{Математическая запись} \\
\midrule
«$A$ больше $B$ на $d$» & $A=B+d$ \\
«$A$ больше $B$ в $k$ раз» & $A=kB$, где $k>1$ \\
«$A$ составляет $\frac pq$ от $B$» & $A=\frac pq B$ \\
«Вместе $S$» & $A+B=S$ \\
\bottomrule
\end{tabular}
\end{center}
\key{Маршрут:} выберите неизвестную величину $\rightarrow$ выразите остальные величины через неё $\rightarrow$ составьте уравнение $\rightarrow$ решите $\rightarrow$ ответьте на вопрос задачи.

\ex{Пример 10. Задача о возрасте}
Мать старше дочери в $4\frac4{11}$ раза, или на $37$ лет. Сколько лет каждой?

Пусть возраст дочери --- $y$ лет. Тогда возраст матери можно записать двумя способами:
\[
\underbrace{4\frac4{11}y}_{\text{в $4\frac4{11}$ раза}}
\qquad\text{и}\qquad
\underbrace{y+37}_{\text{на $37$ лет}}.
\]
Это одно и то же число, поэтому
\[
\frac{48}{11}y=y+37
\quad\Longrightarrow\quad
\left(\frac{48}{11}-\frac{11}{11}\right)y=37
\quad\Longrightarrow\quad
\frac{37}{11}y=37.
\]
\[
y=11,\qquad 11+37=48.
\]
\textcolor{darkaccent}{\textbf{Ответ:}} дочери $11$ лет, матери $48$ лет. Проверка: $48:11=4\frac4{11}$.

\ex{Пример 11. Два арбуза}
Общая масса двух арбузов $13\frac34$ кг. Один составляет $\frac47$ массы другого. Найдите массы.

Пусть масса большего арбуза --- $x$ кг, тогда меньшего --- $\frac47x$ кг. Их общая масса:
\[
x+\frac47x=\frac{55}{4}
\quad\Longrightarrow\quad
\frac{11}{7}x=\frac{55}{4}
\quad\Longrightarrow\quad
x=\frac{55}{4}\cdot\frac7{11}=\frac{35}{4}.
\]
Меньший арбуз: $\frac47\cdot\frac{35}{4}=5$ кг.

\textcolor{darkaccent}{\textbf{Ответ:}} $8\frac34$ кг и $5$ кг. Проверка: $8\frac34+5=13\frac34$.

\clearpage
\sect{6. Задача на догонку: расстояние, скорость, время}
\key{Базовая формула:} $S=vt$ (расстояние равно скорости, умноженной на время). Если один участник стартовал раньше и успел уехать вперёд, это начальное расстояние надо учесть.

\ex{Пример 12. Велосипедист и мотоциклист}
Когда велосипедист находился в $25\frac56$ км от лагеря, вслед за ним выехал мотоциклист и догнал его через $\frac23$ ч. Скорость велосипедиста составляла $\frac38$ скорости мотоциклиста. Найдите обе скорости.

Пусть скорость мотоциклиста $v$ км/ч, тогда скорость велосипедиста $\frac38v$ км/ч.

\begin{center}
\begin{tikzpicture}[x=1cm,y=1cm,>=Stealth]
\draw[darkaccent,thick] (0,0)--(11.1,0);
\draw[accent,line width=1.1pt] (0,-.1)--(0,.22) (5.3,-.1)--(5.3,.22) (10.5,-.1)--(10.5,.22);
\draw[<->,accent,thin] (0,.65)--(5.3,.65);
\node[above,align=center,font=\small] at (2.65,.65) {$25\frac56$ км};
\draw[<->,darkaccent,thin] (5.3,-.5)--(10.5,-.5);
\node[below,align=center,font=\small] at (7.9,-.5) {$\frac38 v\cdot\frac23$};
\draw[<->,muted,thin] (0,-1.35)--(10.5,-1.35);
\node[below,align=center,font=\small] at (5.25,-1.35) {$v\cdot\frac23$};
\node[below,font=\footnotesize] at (0,-.13) {лагерь};
\node[above,font=\footnotesize] at (10.5,.03) {встреча};
\end{tikzpicture}
\end{center}
\vspace{-1mm}
\key{Смысл уравнения:} начальный отрыв плюс путь велосипедиста за $\frac23$ часа равен пути мотоциклиста за то же время.
\[
25\frac56+\frac38v\cdot\frac23=v\cdot\frac23.
\]
\[
\frac{155}{6}+\frac14v=\frac23v
\quad\Longrightarrow\quad
\frac{155}{6}=\left(\frac23-\frac14\right)v=\frac5{12}v.
\]
\[
v=\frac{155}{6}:\frac5{12}=62\ \text{км/ч},
\qquad \frac38v=\frac{93}{4}=23\frac14\ \text{км/ч}.
\]
\textcolor{darkaccent}{\textbf{Ответ:}} мотоциклист --- $62$ км/ч, велосипедист --- $23\frac14$ км/ч.

\clearpage
\sect{7. Что проверить перед записью ответа}
\begin{itemize}
\item Верно ли расставлены скобки и соблюдён порядок действий?
\item Умножили ли Вы на \emph{обратную} дробь при делении? Не оказался ли делитель равным нулю?
\item Сложены ли только числовые коэффициенты у подобных слагаемых? Учтено ли, что $x=1x$?
\item Если в уравнении есть скобка, можно ли сначала найти её значение?
\item При переносе слагаемого через знак равенства изменён ли знак?
\item Не остановились ли Вы на уравнении вида $ax=c$, не найдя $x$?
\item В текстовой задаче различены ли выражения «на столько-то» и «во столько-то раз»?
\item Подставлен ли ответ в исходное уравнение? Указаны ли единицы измерения?
\end{itemize}

\section*{Итоговая памятка}
\begin{center}
\begin{tabularx}{\linewidth}{@{}p{.34\linewidth}X@{}}
\toprule
\textbf{Что видите} & \textbf{Что делаете} \\
\midrule
Деление на дробь & Умножаете на обратную дробь. \\
$a\cdot x=c$ & Находите $x=c:a$. \\
$a:x=c$ & Находите $x=a:c$. \\
$ax+bx$ & Складываете коэффициенты: $(a+b)x$. \\
$a(\ldots)=c$ & Сначала считаете скобку единым множителем. \\
«На $d$ больше» & Составляете сумму $B+d$. \\
«В $k$ раз больше» & Составляете произведение $kB$. \\
Задача на догонку & Складываете начальный отрыв и путь первого участника. \\
\bottomrule
\end{tabularx}
\end{center}
\clearpage

\sect{Тренировочный блок · Путь А}
\textcolor{muted}{10 заданий от базовых действий к составным уравнениям и текстовым задачам. Решайте без калькулятора, записывая основные шаги.}
\begin{enumerate}[label=\textbf{\arabic*.}]
\item Вычислите:
\[
\left(\left(1\frac14\right)^2-0{,}8125\right):0{,}25.
\]
\item Вычислите:
\[
8:0{,}16-12\frac12\cdot0{,}64.
\]
\item Решите уравнение:
\[\frac35x=2\frac1{10}.\]
\item Решите уравнение:
\[\frac57:y=\frac38:3.\]
\item Решите уравнение:
\[\frac29z+\frac59z=2\frac13.\]
\item Решите уравнение:
\[\frac78t+t-\frac12t=2\frac34.\]
\item Решите уравнение:
\[1\frac12\left(\frac23n+\frac16\right)=2\frac14.\]
\item Мать старше сына в $3\frac12$ раза, или на $30$ лет. Сколько лет матери и сыну?
\item Общая масса двух дынь --- $18$ кг. Масса одной составляет $\frac27$ массы другой. Найдите массы обеих дынь.
\item Велосипедист успел отъехать от посёлка на $18$ км, когда за ним отправился мотоциклист. Мотоциклист догнал его через $\frac34$ ч. Скорость велосипедиста составляет $\frac25$ скорости мотоциклиста. Найдите обе скорости.
\end{enumerate}
\clearpage

\sect{Ответы к тренировочному блоку}
\begin{center}
\renewcommand{\arraystretch}{1.42}
\begin{tabularx}{\linewidth}{@{}p{.10\linewidth}X@{}}
\toprule
\textbf{№} & \textbf{Ответ} \\
\midrule
1 & $3$ \\
2 & $42$ \\
3 & $x=3\frac12$ \\
4 & $y=5\frac57$ \\
5 & $z=3$ \\
6 & $t=2$ \\
7 & $n=2$ \\
8 & сыну $12$ лет, матери $42$ года \\
9 & $4$ кг и $14$ кг \\
10 & мотоциклист --- $40$ км/ч, велосипедист --- $16$ км/ч \\
\bottomrule
\end{tabularx}
\end{center}
\vspace{14mm}
\textcolor{darkaccent}{\textbf{Самопроверка:}} совпадение ответа полезно, но особенно важно проверить каждый шаг решения. Для текстовых задач проверьте оба условия: и отношение величин, и их сумму, разность или пройденный путь.
\end{document}
